\chapter{Components and Pages}

\section{Component Tree}

A realistic component tree for the dashboard screen:

\begin{diagrambox}[Task Manager Component Tree]
\begin{verbatim}
App
|-- Navbar
|-- MainLayout
|   |-- Sidebar
|   `-- <Outlet/>
|       |-- Dashboard (page)
|       |   |-- StatsCard x3 (Total/Pending/Done)
|       |   `-- TaskList -> TaskCard (one per task)
|       `-- Profile (page)
\end{verbatim}
\end{diagrambox}

\textbf{Reusable components} (\texttt{StatsCard}, \texttt{TaskCard}) receive data via
props and render it. \textbf{Page components} (\texttt{Dashboard}) own data
fetching. \textbf{Layouts} provide the persistent chrome around the active page.

\section{A Reusable Component}

\begin{lstlisting}[language=JavaScript]
// components/TaskCard.jsx
export default function TaskCard({ task, onStatusChange }) {
  const statusColor = {
    pending: "gray",
    "in-progress": "orange",
    done: "green",
  }[task.status];

  return (
    <div className="task-card">
      <h3>{task.title}</h3>
      <p>{task.description}</p>
      <span className={`badge badge-${statusColor}`}>{task.status}</span>

      {task.status !== "done" && (
        <button onClick={() => onStatusChange(task._id, "done")}>
          Mark Done
        </button>
      )}
    </div>
  );
}
\end{lstlisting}

\texttt{TaskCard} receives \texttt{task} and \texttt{onStatusChange} as props and
never calls the API itself, keeping it reusable and testable.

\section{A Page Component}

\begin{lstlisting}[language=JavaScript]
// pages/Dashboard.jsx
import { useEffect, useState } from "react";
import TaskList from "../components/TaskList";
import StatsCard from "../components/StatsCard";
import { getTasks, updateTaskStatus } from "../services/taskService";

export default function Dashboard() {
  const [tasks, setTasks] = useState([]);
  const [loading, setLoading] = useState(true);

  useEffect(() => {
    getTasks()
      .then((data) => setTasks(data))
      .finally(() => setLoading(false));
  }, []);

  const handleStatusChange = async (id, status) => {
    const updated = await updateTaskStatus(id, status);
    setTasks((prev) => prev.map((t) => (t._id === id ? updated : t)));
  };

  if (loading) return <p>Loading tasks...</p>;

  const pending = tasks.filter((t) => t.status === "pending").length;
  const done = tasks.filter((t) => t.status === "done").length;

  return (
    <div className="dashboard">
      <div className="stats-row">
        <StatsCard label="Total" value={tasks.length} />
        <StatsCard label="Pending" value={pending} />
        <StatsCard label="Done" value={done} />
      </div>
      <TaskList tasks={tasks} onStatusChange={handleStatusChange} />
    </div>
  );
}
\end{lstlisting}

\texttt{Dashboard} owns state and fetching, then hands data down as props --
the "smart page, dumb component" split used throughout the book.

\section{Forms and Conditional Rendering}

\begin{lstlisting}[language=JavaScript]
// components/TaskForm.jsx
import { useState } from "react";

export default function TaskForm({ onCreate }) {
  const [form, setForm] = useState({ title: "", description: "" });
  const [error, setError] = useState("");

  const handleChange = (e) =>
    setForm({ ...form, [e.target.name]: e.target.value });

  const handleSubmit = async (e) => {
    e.preventDefault();
    if (!form.title.trim()) return setError("Title is required");

    await onCreate(form);
    setForm({ title: "", description: "" });
    setError("");
  };

  return (
    <form onSubmit={handleSubmit}>
      <input name="title" value={form.title} onChange={handleChange} />
      <textarea
        name="description"
        value={form.description}
        onChange={handleChange}
      />
      {error && <p className="error">{error}</p>}
      <button type="submit">Create Task</button>
    </form>
  );
}
\end{lstlisting}

\begin{tipbox}[TIP]
The spread operator updates one field without losing others; a computed key
(\texttt{[e.target.name]}) lets one \texttt{handleChange} serve every input.
\end{tipbox}
